% !TeX root = ../main.tex
% ----------------------------------------------------------------------
% Encabezados, títulos, índice y enlaces
% ----------------------------------------------------------------------
\usepackage{epigraph}
\usepackage{fancyhdr}
\usepackage{titlesec}
\usepackage{titletoc}
\usepackage{emptypage}

\addto\captionsspanish{%
  \renewcommand{\contentsname}{Contenido}%
  \renewcommand{\bibname}{Bibliografía}%
}

% --- Epígrafe del prefacio ---
\renewcommand{\epigraphflush}{flushright}
\renewcommand{\epigraphsize}{\normalsize\itshape}
\setlength{\epigraphwidth}{0.7\textwidth}
\renewcommand{\textflush}{flushepinormal}
\renewcommand{\sourceflush}{flushepinormal}

% --- Encabezado y pie de página ---
\newcommand{\fechaclase}{}
\newcommand{\clase}[2]{\renewcommand{\fechaclase}{#1}\chapter{#2}}
\newcommand{\capitulosinnumero}[1]{%
  \chapter*{#1}\addcontentsline{toc}{chapter}{#1}\markboth{#1}{}%
}

\pagestyle{fancy}
\fancyhf{}
\fancyhead[L]{\small\textcolor{darkblue}{\textbf{Notas del teórico}}}
\fancyhead[R]{\small\textcolor{mediumgray}{\nouppercase{\leftmark}}}
\fancyfoot[C]{\thepage}
\renewcommand{\headrulewidth}{0.5pt}
\renewcommand{\headrule}{\color{darkblue}\hrule width\headwidth height\headrulewidth \kern-\headrulewidth}
\renewcommand{\footrulewidth}{0pt}
\renewcommand{\chaptermark}[1]{\markboth{Clase \thechapter\ \textbullet\ #1}{}}

\fancypagestyle{plain}{%
  \fancyhf{}%
  \fancyfoot[C]{\thepage}%
  \renewcommand{\headrulewidth}{0pt}%
}

% --- Títulos ---
\titleformat{\chapter}[display]
  {\normalfont}
  {\large\bfseries\color{accentblue}Clase \thechapter%
    \normalfont\normalsize\color{mediumgray}\quad\textbullet\quad\fechaclase}
  {4pt}
  {\Huge\bfseries\color{darkblue}}
  [\vspace{0.4em}{\color{darkblue}\titlerule[1.5pt]}]
\titleformat{name=\chapter, numberless}[display]
  {\normalfont}
  {}
  {0pt}
  {\Huge\bfseries\color{darkblue}}
  [\vspace{0.4em}{\color{darkblue}\titlerule[1.5pt]}]
\titlespacing*{\chapter}{0pt}{-1.5em}{2em}

\titleformat{\section}
  {\Large\bfseries\color{darkblue}}
  {\thesection.}{0.5em}{}
  [\vspace{0.15em}{\color{darkblue}\titlerule[0.6pt]}]
\titleformat{\subsection}
  {\large\bfseries\color{accentblue}}
  {\thesubsection.}{0.5em}{}
\titleformat{\subsubsection}
  {\normalsize\bfseries\color{accentblue}}
  {\thesubsubsection.}{0.5em}{}
\titlespacing*{\section}{0pt}{2em plus 0.5em minus 0.2em}{0.8em}
\titlespacing*{\subsection}{0pt}{1.5em plus 0.4em minus 0.2em}{0.4em}

% --- Índice general ---
\setcounter{tocdepth}{1}
\titlecontents{chapter}[5.2em]
  {\addvspace{0.9em}\bfseries\color{darkblue}}
  {\contentslabel[Clase~\thecontentslabel]{5.2em}}
  {\hspace*{-5.2em}}
  {\hfill\contentspage}
\titlecontents{section}[7.7em]
  {\small}
  {\contentslabel{2.5em}}
  {\hspace*{-2.5em}}
  {\titlerule*[0.6pc]{.}\contentspage}

% --- Enlaces ---
\usepackage{hyperref}
\hypersetup{
  colorlinks=true,
  linkcolor=darkblue,
  urlcolor=darkblue,
  citecolor=darkblue,
  bookmarksnumbered=true,
  breaklinks=true,
  pdftitle={Notas del teórico - Complementos de Análisis Matemático},
  pdfauthor={Jordi Bustos}
}

% --- Anclas únicas para los entornos que comparten contador ---
\def\theHdefinition{\thechapter.\arabic{definition}}
\def\theHprop{\thechapter.\arabic{definition}}
\def\theHtheorem{\thechapter.\arabic{definition}}
\def\theHlemma{\thechapter.\arabic{definition}}
\def\theHcorollary{\thechapter.\arabic{definition}}

